\subsection{自由模中一个元素的内容}

设 A 为主理想整环，L 为自由 A-模，x 为 L 的一个元素。当 f 遍历 L 上线性形式的集合 \(L^{*}\) 时，诸元素 \(f(x)\) 构成 A 的一个理想 \(\mathfrak{c}_{\mathrm{L}}(x)\)，称为 x 在 L 中的内容。A 的元素 c 称为 x 在 L 中的一个内容，如果它生成理想 \(\mathfrak{c}_{\mathrm{L}}(x)\)；这相当于说，存在 L 上的一个线性形式 f，使得 \(f(x)=c\)，并且 c 整除 L 上每个线性形式 g 所对应的 \(g(x)\)。设 \((e_{i})_{i\in I}\) 为 L 的一组基；令 \(x=\sum a_{i}e_{i}, a_{i}\in A\)；则理想 \(\mathfrak{c}_{\mathrm{L}}(x)\) 由诸和式 \(\sum a_{i}b_{i}\) 组成，其中 \((b_{i})\) 遍历集合 \(A^{I}\)；由此立即得出，A 的元素 c 是 x 在 L 中的一个内容，当且仅当它是 x 的坐标族 \((a_{i})\) 的最大公因子。

我们称 x 是不可除元，如果 \(\mathfrak{c}_{L}(x)=\mathbf{A}\)，即 x 关于 L 的一组基的坐标整体互素。

引理1. --- 设L是主理想整环A上的自由模，且x是L的一个元素。则以下条件等价：

\begin{enumerate}
\def\labelenumi{(\roman{enumi})}
\item
  \(x\) 是不可除的；
\item
  存在一个线性形式 \(f\)，它定义在 \(L\) 上，使得 \(f(x) = 1\)；
\item
  \(x\) 非零，且子模 \(\mathbf{A}x\) 是 \(\mathbf{L}\) 的一个直因子;
\item
  \(x\) 是 \(\mathbf{L}\) 的某组基中的一个元素.
\item
  \(\Rightarrow\) (ii)：这由定义得出。
\item
  \(\Rightarrow\) (iii)：设 \(f\) 为 \(\mathbf{L}\) 上的一个线性形式，使得 \(f(x) = 1\)；则 \(x \neq 0\)，且映射 \(y \mapsto f(y)x\) 是 \(\mathbf{L}\) 的一个投影算子，其像为 \(\mathbf{A}x\).
\item
  \(\Rightarrow\) (iv)：设 \(L'\) 是 \(Ax\) 在 \(L\) 中的一个补，并设 \(B'\) 是 \(L'\) 的一组基（VII，第15页，推论2）；则 \(B' \cup \{x\}\) 是 \(L\) 的一组基.
\item
  \(\Rightarrow\) (i)：平凡。
\end{enumerate}

注记。--- 1) 若 \(x\) 是 \(L\) 中的非零元素，且 \(c\) 是 \(x\) 的一个内容，则存在唯一的元素 \(y\)（属于 \(L\)）使得 \(x = cy\)；把这个元素记作 \(x / c\)；则 \(x / c\) 是 \(L\) 的一个不可除元素。

\begin{enumerate}
\def\labelenumi{\arabic{enumi})}
\setcounter{enumi}{1}
\tightlist
\item
  内容 \(\mathfrak{c}_{\mathrm{L}}(x)\) 是挠模 \(\mathrm{L} / \mathrm{A}x\) 的零化子.
\end{enumerate}

设 L 是主理想整环 A 上的自由模，且 M 是 L 的一个子模；由 VII 第 4 页引理 1，族 \(\mathfrak{c}_{\mathrm{L}}(x)\)（\(x \in M\)）具有极大元；若 \(M \neq \{0\}\)，则这样的极大元是非零的。

命题3. --- 设L是主理想整环A上的自由模，且M是L的一个非零子模。设x是M中使 \(\mathfrak{c}_{\mathrm{L}}(x)\) 在M的各元素的内容之中为极大的一个元素，设c是x在L中的一个内容，并设f是L上的一个线性形式，使得 \(f(x) = c\).

\begin{enumerate}
\def\labelenumi{\alph{enumi})}
\item
  L 是 A(x/c) 与 f 的核 K 的直和。
\item
  \(\mathbf{M}\) 是 \(\mathbf{A}\mathbf{x}\) 与 \(\mathbf{K} \cap \mathbf{M}\) 的直和.
\item
  \(g(M) \subset Ac\) 对每个线性形式 \(g\)（\(L\) 上的）都成立。
\end{enumerate}

令 \(y = x / c\)；显然 \(\mathbf{A}y \cap \mathbf{K} = \{0\}\)，因为 \(f(y) = 1\)。此外，对每个 \(u \in L\)，我们有

\[
u = f (u) y + (u - f (u) y),
\]

其中 \(f(u) y \in \mathbf{A}y\)，且 \(u - f(u)y \in \mathbf{K}\)；这就证明了 \(a\)）。现在注意，对 \(u \in \mathbf{M}\) 我们有 \(f(u) \in \mathbf{A}c\)：事实上，设 \(u \in \mathbf{M}\)，并设 \(d\) 是 \(f(u)\) 与 \(c\) 的一个最大公因子；则存在 \(\lambda, \mu \in \mathbf{A}\) 使得 \(d = \lambda f(u) + \mu c = f(\lambda u + \mu x)\)；因而元素 \(\lambda u + \mu x\)（\(\mathbf{M}\) 中的）的内容整除 \(d\)；由 \(c\) 的极大性，这蕴含 \(d\) 与 \(c\) 相伴，从而 \(f(u) \in \mathbf{A}c\)。于是对每个 \(u\)（在 \(\mathbf{M}\) 中），我们可以写出

\[
u = (f (u) / c) x + (u - (f (u) / c) x) \in \mathrm{A} x + (\mathrm{K} \cap \mathrm{M}),
\]

这表明 \(b\)）。最后，设 \(g\) 是 \(L\) 上的一个线性形式；由 \(a\)）存在一个标量 \(\alpha \in A\) 和一个线性形式 \(h\)（\(K\) 上的），使得 \(g(u) = \alpha f(u) + h(u - f(u)y)\)；于是由 b) 我们有 \(g(\mathbf{M}) \subset \mathbf{A}c + h(\mathbf{K} \cap \mathbf{M})\)。因此要证明 c)，只需证明对 K 上的每个线性形式 h 都有 \(h(\mathbf{K} \cap \mathbf{M}) \subset \mathbf{A}c\)，或等价地，对 L 上满足 \(h(x) = 0\) 的每个线性形式 h 成立；现在，若 \(u \in K \cap M\) 且 d 是 \(h(u)\) 与 c 的一个最大公因子，则存在 \(\lambda\)、\(\mu \in A\) 使得 \(d = \lambda h(u) + \mu c\)；则 \((f + h)(\lambda u + \mu x) = d\)，这如上所述蕴含 \(h(u) \in Ac\)，从而 c) 成立。

